\section{Recurrent singular acquisition on accumulating components}\label{sec:moran}
We verify the preceding theorem from a raw kernel whose acquired measures need not have a quantization dimension. This verification uses no covariance, differentiable chart, or stationary environment resolvent. The conditional laws are explicit singular product measures; the component masses evolve by the raw refresh flow.

\subsection{A nonhomogeneous binary construction}
Let $(\rho_n)_{n\ge1}$ be any sequence in $[1/5,1/3]$, put $\ell_0=1$ and $\ell_n=\prod_{i=1}^n\rho_i$, and define
\begin{equation}\label{eq:moran-map}
 x(\omega)=\sum_{n\ge1}(1-\rho_n)\ell_{n-1}\omega_n,
 \qquad \omega\in\{0,1\}^{\N}.
\end{equation}
The series is uniformly convergent. Its image $C$ is compact and its digit representation is unique, because the two children at step $n$ are separated by the gap $(1-2\rho_n)\ell_{n-1}>0$. Let $\sigma$ be the image of independent fair digits. A cylinder of length $n$ has mass $2^{-n}$ and diameter $\ell_n$; the endpoints with constant tails belong to $C$.

For $k\ge1$ set
\begin{equation}\label{eq:moran-profile}
 r(k)=\ell_{\lfloor\log_2k\rfloor},\qquad r(0)=1.
\end{equation}
No assumption that $-n^{-1}\log\ell_n$ converges is imposed.

\begin{lemma}[Mass and cover at every binary scale]\label{lem:moran-profile}
There are at most $k$ points of $C$ covering $C$ with radius $r(k)$. For every real center $z$,
\[
 \sigma(B(z,r(k)/100))\le(2k)^{-1}.
\]
Moreover, $r(k-1)^2\le25r(k)^2$ for $k\ge2$. The digit shift
$F(x(\omega))=x(\omega_2,\omega_3,\ldots)$ preserves $\sigma$ and is $15$-Lipschitz on $C$.
\end{lemma}
\begin{proof}
Put $n=\lfloor\log_2k\rfloor$. Selecting one valid point from each length-$n$ cylinder uses $2^n\le k$ points and has covering radius at most $\ell_n$. Distinct length-$(n+2)$ cylinders have mutual distance at least
\[
 \min_{1\le i\le n+2}(1-2\rho_i)\ell_{i-1}
       \ge\ell_{n+1}/3\ge\ell_n/15.
\]
A ball of diameter $\ell_n/50$ therefore meets at most one such cylinder, whose mass is $2^{-(n+2)}\le(2k)^{-1}$. This holds for centers outside $C$ as well. The regularity bound changes only when $k$ is a power of two, when it follows from $\rho_n\ge1/5$.

Shift invariance follows because the shifted digits are again independent and fair, although the coding coefficients are nonhomogeneous. If two sequences first differ at digit $i\ge2$, their original distance is at least $(1-2\rho_i)\ell_{i-1}$, whereas their shifted distance is at most $\ell_{i-2}$. Their ratio is at most $1/[(1-2\rho_i)\rho_{i-1}]\le15$. If $i=1$, use shifted diameter at most one and original distance at least $1-2\rho_1\ge1/3$. This proves the global Lipschitz bound.
\end{proof}

\subsection{The prepared experiment}
Index components by $j\ge0$ and set $a_j=10^{-(j+2)}$. The hidden physical state is a prediction probability
\[
 p=\tfrac12+a_j(4+x),\qquad x\in C,
\]
or the limiting anchor $o=1/2$. Let $D_j=\tfrac12+a_j(4+C)$ and let $\sigma_j$ be the corresponding image of $\sigma$. The union of the $D_j$ and $o$ is compact in the ordinary absolute-value metric. The sets $D_j$ lie in disjoint intervals of length $a_j$ accumulating at $o$; their $a_j/2$ neighborhoods are disjoint. The task is the single executable terminal probe $T\sim\operatorname{Bernoulli}(p)$, with square loss. This probe is generated only after the prediction has been committed.

One paid initialization draws $j$ with known law $w_0$, draws $x$ with law $\sigma$, and reveals the resulting state. At continuation step $t$, a report type $H,E,R$ is chosen with probabilities $h_t,u_t,v_t$, independently of the current state. On $H$ the state is unchanged and only the type is reported. On $E$ the state becomes $1/2+a_j(4+F(x))$ and only the type is reported. On $R$ a new component is drawn with law $\pi_t$, a new independent point with law $\sigma_j$ is drawn, and the new state and type are reported. The anchor is fixed by $H$ and $E$; $R$ has the same state-independent law there. Every call, including a hold, costs one raw acquisition and one unit of physical time. These are nonnegative normalized kernels on compact standard-Borel spaces. A reveal can equivalently report its uniquely determined component and digit sequence; no hidden component is thereby supplied to the terminal decoder after the memory cut.

A finite protocol has $n$ continuation calls after the paid initialization, so its preforecast raw-call count is $N=n+1$. The terminal probe is a further paid call. Nonstationary probabilities are known exogenous protocol data. There are no optional control actions in this realization beyond continue and the declared terminal command; controller optimization is not smuggled into a fixed-protocol statement.

The observer's exact predictive state is $p$. Indeed the initial reveal and every later update determine $p$ from the full history. Equal $p$ values give the same future laws by the explicit kernel; different $p$ values are separated by the terminal Bernoulli probe. Thus the quotient is exactly the displayed compact set with the declared protocol interface. The shift is continuous on each component and at the anchor, because component diameters tend to zero. Reports $H,E$ have state-independent probabilities, and the entire reveal law under $R$ is state independent. Report total variation between any two preterminal states is consequently zero; terminal Bernoulli total variation equals $|p-p'|$. This supplies report continuity without falsely treating Dirac reveal laws as globally TV-continuous in an unrevealed state. The reveal is a freshly drawn state, not a noiseless observation of an arbitrary old state.

\begin{theorem}[Nonstationary recurrent singular realization]\label{thm:moran}
Fix a contraction sequence in $[1/5,1/3]$, any prepared law $w_0$ on the countable components, and any known nonstationary raw-flow protocol satisfying \eqref{eq:flow-balance} with $G=15$, some $\eta>0$, and $K>(1+\eta)225-1$. Let $w_n$ be given by \eqref{eq:raw-flow}, and define
\begin{equation}\label{eq:moran-Q}
 \mathcal Q_n(M)=\inf_{\sum_jk_j\le M}
       \sum_{j\ge0}w_{n,j}a_j^2
       \ell_{\lfloor\log_2(\max\{1,k_j\})\rfloor}^{\,2}.
\end{equation}
For every $n\ge0$ and $M\ge2$,
\begin{equation}\label{eq:moran-risk}
 R_\cp(n+1,M)-B_{n+1}\asymp
 R_\on(n+1,M)-B_{n+1}\asymp\mathcal Q_n(M).
\end{equation}
The constants depend on $K,\eta$ and the fixed scale-separation constants, but not on $n$, the number or rarity of acquired components, or the holding probabilities. The same statement applies to a finite subset of components with arbitrary positive scales obeying the displayed relative separation. No power-law or limiting dimension for $\sigma$ is required.
\end{theorem}
\begin{proof}
By Lemma~\ref{lem:moran-profile}, \textup{(MG)} holds with $r_j(k)=a_jr(k)$, $C_g=5$, $c_g=1/100$, $c_s=1/2$, and $D_r=25$; the factor five allows discarding to the anchor. The raw computation in Proposition~\ref{prop:recharge-flow} gives
$\mu_n=\sum_jw_{n,j}\sigma_j$. These are the acquired laws, not an auxiliary choice of measure on a reachable set.

Here is the local machine for an allocation. On an allocated component with $k_j>0$, retain its component index and a length-$\lfloor\log_2k_j\rfloor$ prefix, represented by the point with that prefix and all-zero tail. The entire pair is one of at most $\sum_j k_j$ labels. A reset into that component reads the prefix and overwrites the label. A reset into an unallocated component writes the common anchor label. A hold copies the label exactly. On $E$, an allocated label shifts its stored digit word and appends zero; the anchor remains the anchor. Thus the machine never needs to recover the index of an unallocated component.

On an allocated component, shifting the representative and rounding within the same component has error at most $15e_{t-1}+a_jr(k_j)$; indeed the shifted all-zero-tail representative is already a valid grid point, so the latter term is a harmless upper allowance. On an unallocated component, the distance to the anchor is at most $5a_j$. Under $E$, the ratio of the new distance to the old is at most $5/4$, so $15e_{t-1}$ still bounds it. On reset the old error vanishes and the new error is at most $5a_jr(k_j)$, including the zero-allocation case. Thus \eqref{eq:bridge-local} holds with $R_j=5a_jr(k_j)$ and $b_t\le1$. It is executable with the stated labels, not with an uncharged mode register. The initialized error measure has normalized bound one. Weighted continuation invariance follows from shift invariance and its gain bound. Proposition~\ref{prop:recharge-flow} now verifies \textup{(TB)} from the raw flow condition, and Theorem~\ref{thm:bridge} proves \eqref{eq:moran-risk}. The initialization/continuation indexing is exactly $N=n+1$.
\end{proof}

The class contains repeated active expansions. For example, with $\eta=1$, take $K=500$, $u_t=\theta_t/2001$, $v_t=2000\theta_t/2001$, $h_t=1-\theta_t$, and $\pi_t=(w_{t-1}+\omega_t)/2$. Then \eqref{eq:flow-balance} holds with room to spare, and arbitrary sequences $0\le\theta_t\le1$ are allowed. Whenever active expansions occur infinitely often they are recurrent, not a single initial transient. Arbitrarily long intervals with $\theta_t=0$ add raw time but no new rounding. Nonstationarity of $\omega_t$ can move acquired mass among old and newly introduced components. No constant divides by $\theta_t$ or by $w_{n,j}$.

\subsection{An oscillatory singular law and finite descriptions}
\begin{corollary}[A sharp causal profile without a limiting dimension]\label{cor:no-dimension}
There is a computably specified choice of $(\rho_n)$ for which the one-component squared quantization error $q_M(\sigma)$ satisfies
\[
 q_M(\sigma)\asymp\ell_{\lfloor\log_2M\rfloor}^{\,2},
\]
and $\log q_M(\sigma)/\log M$ has no limit. The checkpoint and online excess risks in the one-component version of Theorem~\ref{thm:moran} have the same nonconvergent profile, with uniform constants through all scales.
\end{corollary}
\begin{proof}
Use alternating blocks of $\rho_n=1/3$ and $\rho_n=1/5$, whose lengths are $2^{2^k}$ for block index $k\ge1$. Each new block length dominates the sum of all previous lengths. The cover and small-ball proof of Theorem~\ref{thm:bridge}, with a single component, gives the displayed estimate for every $M$ (the anchor is immaterial up to the one-label inequality). At $M=2^n$ the quotient of logarithms is
$2\sum_{i\le n}\log\rho_i/(n\log2)+o(1)$.
Along alternate block endpoints it tends to $-2\log3/\log2$ and $-2\log5/\log2$. Uniform multiplicative constants do not change these limits. Theorem~\ref{thm:moran} transfers the same profile to the online class.
\end{proof}

A finite-bit version has a concrete interpretation. Choose a finite active set of component indices and prefix lengths $d_j=\lfloor\log_2k_j\rfloor$. A reset packet supplies the component index and the first $d_j$ digits on those components; all other indices need only be mapped to the anchor by a declared input interface. This truncation is itself the observation kernel of the finite-input experiment, not free access to an infinite word. The representative states are integer labels. Shifts and holds are exact integer operations; a terminal probability is evaluated from a finite sum in \eqref{eq:moran-map}, followed by a declared output rounding. If $\rho_n$ and $a_j$ are rational with explicitly bounded bit length through the used depth, the operation counts and workspace are finite and computable from those lengths and the maximum depth. Surviving digits belong to the persistent label; temporary evaluation words are erased before the next call.

For a countable law, a certified tail bound $\sum_{j>J}w_{n,j}a_j^2$ and computable prefix weights permit finite enumeration of allocations up to any prescribed additive profile tolerance. A constant-factor allocation additionally needs a certified positive lower scale. The general theorem does not turn an arbitrary real sequence or an unbounded integer report into a finite-cost routine. Even without such effectiveness data, its Borel finite-label assertion remains precise. This separation is why program and input-interface assumptions are retained in the resource ledger.
